\chapter{Jaminan Kualitas Data: Pengujian Otomatis dan Rekonsiliasi}
\label{modul:jaminan-kualitas-data}

\section{Identitas Modul}

\begin{identitasmodul}
\textbf{Sumber materi}    & Pertemuan 13 \\
\textbf{CPMK}             & CPMK-102 \\
\textbf{Minggu pelaksanaan} & Minggu ke-13, setelah kuliah Pertemuan 13 \\
\textbf{Durasi tatap muka} & 120 menit \\
\textbf{Estimasi tugas}   & 4--5 jam di luar sesi \\
\textbf{Moda kerja}       & Individual \\
\textbf{Perangkat}        & PostgreSQL 17 dan dbt Core \\
\textbf{Dataset}          & Satu batch NusaMart yang memuat anomali terkendali \\
\textbf{Skrip pendukung}  & \berkas{sql/modul-10/} beserta \berkas{dbt/modul-10/} \\
\textbf{Luaran}           & Skrip profiling dan rekonsiliasi, minimal delapan uji data, serta laporan kualitas data per batch \\
\end{identitasmodul}

\slideref{13}

\section{Capaian dan Indikator Keberhasilan}

Mahasiswa mampu menerjemahkan aturan bisnis menjadi pemeriksaan kualitas yang
terukur, mendeteksi anomali secara otomatis, menangani data gagal tanpa mengubah
data mentah, dan membuktikan kesesuaian angka antara sumber dan warehouse.

\begin{capaian}
Pada akhir modul ini mahasiswa mampu:
\begin{enumerate}[leftmargin=1.4em]
  \item memprofilkan satu batch dan menemukan anomali yang belum diketahui
        sebelumnya;
  \item menerjemahkan aturan bisnis menjadi generic dan custom data test yang
        dapat dijalankan otomatis;
  \item menetapkan tingkat keparahan dan ambang setiap uji, serta menjelaskan
        alasannya;
  \item mengarantina baris yang gagal tanpa mengubah data mentah maupun
        menghapusnya;
  \item merekonsiliasi empat besaran antara sumber dan gudang, dan menjelaskan
        setiap selisih yang muncul.
\end{enumerate}
\end{capaian}

Sembilan modul sebelumnya mengandaikan angka yang dimuat adalah benar. Modul ini
membalik andaian itu: data diperlakukan sebagai sesuatu yang harus
\emph{dibuktikan} layak dipercaya, dan pembuktian itu harus dapat dijalankan
ulang setiap kali data baru masuk.

\section{Ruang Lingkup}

Modul ini melanjutkan pengenalan generic data tests pada Modul 8. Fokusnya
adalah perancangan aturan kualitas, custom test, threshold, karantina data
gagal, serta rekonsiliasi jumlah baris dan control total per batch. Mahasiswa
menilai completeness, uniqueness, validity, consistency, integrity, dan
timeliness pada pipeline NusaMart.

\section{Prasyarat dan Persiapan}

\subsection{Prasyarat}

\begin{itemize}
  \item Proyek dbt yang berjalan dari Modul 8, termasuk empat generic test;
  \item pipeline ETL dan tabel \perintah{audit\_muat} dari Modul 7;
  \item koneksi \perintah{postgres\_fdw} ke ketiga sumber --- rekonsiliasi
        lintas basis data baru mungkin sejak modul itu;
  \item hasil profiling Modul 1, yang akan dibandingkan dengan profiling hari
        ini;
  \item pengertian baris unknown dan constraint dari Modul 5.
\end{itemize}

\subsection{Pre-lab}

\begin{enumerate}[leftmargin=1.4em]
  \item Muat batch beranomali yang disediakan asisten:
\begin{lstlisting}[style=shellgaya]
docker compose exec -T postgres_source \
  psql -U praktikum -d nusamart_oltp -f /seed/batch_anomali.sql
\end{lstlisting}
  \item Jalankan pipeline Modul 7 untuk batch tersebut, lalu \perintah{dbt build}.
        Perhatikan: seluruhnya akan berjalan \textbf{tanpa galat}. Itulah
        persoalannya.
  \item Buka dashboard Modul 9 dan catat angka totalnya. Angka itu tampak
        wajar --- dan sebagian di antaranya salah.
  \item Siapkan jawaban berikut pada \berkas{modul-10-kualitas/pre-lab.md}:
  \begin{enumerate}[label=\alph*.,leftmargin=1.4em]
    \item Constraint basis data Modul 5 sudah menolak kuantitas tidak positif.
          Sebutkan satu jenis kesalahan data yang \emph{tidak} dapat ditolak
          oleh constraint mana pun.
    \item Sebuah uji melaporkan 12 baris gagal dari 5 juta. Apakah pemuatan
          harus dibatalkan? Apa yang menentukan jawabannya?
    \item Bila ditemukan baris yang salah, mengapa memperbaikinya langsung di
          basis data sumber adalah tindakan yang keliru?
  \end{enumerate}
\end{enumerate}

\section{Konsep Dasar}

\subsection{Enam Dimensi Kualitas Data}

``Data berkualitas'' terlalu kabur untuk diuji. Ia harus dipecah menjadi
dimensi yang masing-masing dapat diukur dengan satu angka
\citep{batini2016, dama2017}.

\begin{center}
\small
\begin{tabular}{@{}p{0.18\textwidth}>{\raggedright\arraybackslash}p{0.32\textwidth}>{\raggedright\arraybackslash}p{0.42\textwidth}@{}}
\toprule
\textbf{Dimensi} & \textbf{Pertanyaannya} & \textbf{Contoh ukuran pada NusaMart} \\
\midrule
Completeness & Adakah yang hilang? &
  Persentase baris fakta yang menunjuk baris unknown \\
Uniqueness & Adakah yang kembar? &
  Jumlah \perintah{(toko, nomor\_struk)} yang muncul dua kali \\
Validity & Sesuaikah bentuknya? &
  Jumlah kuantitas yang tidak positif \\
Consistency & Cocokkah antarbagian? &
  Baris yang \perintah{subtotal}-nya tidak sama dengan hitungan \\
Integrity & Utuhkah relasinya? &
  Baris fakta yang produknya tidak ada di dimensi \\
Timeliness & Cukup segarkah? &
  Selisih hari antara transaksi dan pemuatannya \\
\bottomrule
\end{tabular}
\end{center}

\begin{catatan}
Perhatikan bahwa setiap ukuran berupa \emph{angka}, bukan pernyataan
ya-atau-tidak. Angka memungkinkan penetapan ambang, perbandingan antarbatch,
dan penilaian apakah keadaan membaik atau memburuk --- tiga hal yang tidak
mungkin dilakukan dengan penilaian ``bagus'' atau ``jelek''.
\end{catatan}

\subsection{Profiling sebagai Penemuan}

Profiling Modul 1 dikerjakan untuk \emph{memahami} sumber. Profiling di sini
dikerjakan untuk \emph{menemukan} yang tidak diharapkan --- dan keduanya
memakai pemeriksaan yang sama.

Perbedaannya terletak pada pembanding. Sekarang Anda memiliki profil batch-batch
sebelumnya, sehingga yang dicari bukan hanya nilai yang mustahil, melainkan
juga nilai yang \emph{berubah tanpa sebab}.

\begin{konsep}{Tiga jenis anomali}
\begin{enumerate}[leftmargin=1.4em]
  \item \textbf{Mustahil} --- kuantitas negatif, tanggal di masa depan.
        Ditemukan tanpa pembanding.
  \item \textbf{Menyimpang} --- rata-rata harga naik sepuluh kali lipat dalam
        semalam. Hanya ditemukan dengan membandingkan batch.
  \item \textbf{Hilang} --- satu toko yang biasanya mengirim seribu transaksi
        hari ini mengirim nol. Paling sulit ditemukan, karena
        \emph{ketiadaan} tidak menghasilkan baris yang bisa diperiksa.
\end{enumerate}
\end{konsep}

Jenis ketiga inilah alasan rekonsiliasi jumlah baris tidak dapat digantikan
oleh uji apa pun atas isi tabel.

\subsection{Generic Test, Custom Test, dan Ambang}

Uji baku Modul 8 menangkap kerusakan struktural. Aturan bisnis memerlukan uji
yang ditulis sendiri.

\begin{center}
\small
\begin{tabular}{@{}p{0.20\textwidth}>{\raggedright\arraybackslash}p{0.34\textwidth}>{\raggedright\arraybackslash}p{0.38\textwidth}@{}}
\toprule
\textbf{Jenis} & \textbf{Bentuknya} & \textbf{Contoh} \\
\midrule
Generic & Beberapa baris YAML &
  \perintah{not\_null}, \perintah{unique}, \perintah{relationships} \\
Singular & Satu berkas SQL yang mengembalikan baris gagal &
  Satu produk tidak boleh punya dua versi aktif \\
\bottomrule
\end{tabular}
\end{center}

\begin{konsep}{Uji singular pada dbt}
Sebuah uji singular adalah query yang \textbf{mengembalikan baris yang
melanggar aturan}. Uji dinyatakan lulus bila query itu tidak mengembalikan
baris apa pun.

Logika terbalik ini sering membingungkan pada awalnya: yang ditulis bukan
``buktikan aturannya benar'', melainkan ``tunjukkan yang melanggarnya''.
\end{konsep}

Tidak semua kegagalan setara. dbt membedakan dua tingkat keparahan, dan
pembedaan itu yang membuat pengujian dapat dipakai di lingkungan sebenarnya.

\begin{center}
\small
\begin{tabular}{@{}p{0.16\textwidth}>{\raggedright\arraybackslash}p{0.34\textwidth}>{\raggedright\arraybackslash}p{0.44\textwidth}@{}}
\toprule
\textbf{Tingkat} & \textbf{Akibatnya} & \textbf{Dipakai untuk} \\
\midrule
\perintah{error} & Build berhenti &
  Kerusakan yang membuat angka salah: kunci kembar, relasi putus \\
\perintah{warn} & Build lanjut, dilaporkan &
  Keadaan yang perlu diketahui tetapi tidak membatalkan: baris unknown
  meningkat \\
\bottomrule
\end{tabular}
\end{center}

\begin{peringatan}
Menjadikan seluruh uji bertingkat \perintah{error} tampak paling aman, tetapi
justru berbahaya. Pipeline yang berhenti setiap malam karena dua belas baris
bermasalah akan segera dimatikan uji-nya oleh orang yang lelah dibangunkan ---
dan sesudah itu tidak ada lagi yang memeriksa apa pun.

Ambang yang dipilih harus dapat dijelaskan: mengapa 12 baris ditoleransi
sedangkan 1.200 tidak.
\end{peringatan}

\subsection{Karantina Data Gagal}

Ketika sebuah baris gagal uji, ada tiga pilihan, dan hanya satu yang benar.

\begin{center}
\small
\begin{tabular}{@{}p{0.22\textwidth}>{\raggedright\arraybackslash}p{0.30\textwidth}>{\raggedright\arraybackslash}p{0.42\textwidth}@{}}
\toprule
\textbf{Pilihan} & \textbf{Akibatnya} & \textbf{Penilaian} \\
\midrule
Membuang & Baris hilang tanpa jejak &
  Salah: total menyusut tanpa penjelasan \\
Memperbaiki di sumber & Bukti asli lenyap &
  Salah: sumber adalah catatan sistem lain, bukan milik gudang \\
Mengarantina & Baris disimpan terpisah beserta alasannya &
  Benar \\
\bottomrule
\end{tabular}
\end{center}

Karantina menyimpan baris gagal pada tabel tersendiri, lengkap dengan aturan
yang dilanggarnya dan batch asalnya. Barisnya tidak masuk mart --- sehingga
laporan tidak tercemar --- tetapi juga tidak hilang, sehingga jumlahnya dapat
dihitung dan diperiksa manusia.

\begin{catatan}
Aturan ``data mentah tidak diubah'' sudah berlaku sejak Modul 5 sebagai kaidah
reproduksibilitas. Di sini ia memperoleh alasan keduanya: bila data sumber
diperbaiki, rekonsiliasi antara sumber dan gudang akan selalu cocok --- dan
kecocokan yang diperoleh dengan cara itu tidak membuktikan apa pun.
\end{catatan}

\subsection{Rekonsiliasi Sumber dan Gudang}

Rekonsiliasi menjawab satu pertanyaan: apakah yang ada di gudang sama dengan
yang ada di sumber? Empat besaran diperiksa, dan keempatnya diperlukan.

\begin{center}
\small
\begin{tabular}{@{}p{0.26\textwidth}>{\raggedright\arraybackslash}p{0.64\textwidth}@{}}
\toprule
\textbf{Besaran} & \textbf{Menangkap} \\
\midrule
Jumlah baris & Baris yang hilang atau kembar saat pemuatan \\
Transaksi unik & Baris kembar yang tidak menambah jumlah baris karena
  menggantikan yang lain \\
Total kuantitas & Kesalahan pada kolom measure yang tidak mengubah jumlah
  baris \\
Total nilai penjualan & Kesalahan konversi satuan atau pembulatan \\
\bottomrule
\end{tabular}
\end{center}

Keempatnya diperlukan karena masing-masing buta terhadap yang lain. Jumlah baris
yang cocok tidak menjamin nilainya cocok; nilai yang cocok tidak menjamin tidak
ada baris yang tertukar.

\begin{peringatan}
Selisih rekonsiliasi \textbf{tidak harus nol}. Yang wajib adalah setiap selisih
\emph{dapat dijelaskan} --- transaksi batal yang sengaja disaring, baris
terkarantina, atau transaksi terlambat yang belum sampai. Selisih yang
dijelaskan adalah temuan; selisih yang didiamkan adalah kegagalan.
\end{peringatan}

\section{Studi Kasus dan Protokol Praktikum}

Batch yang dimuat pada pre-lab memuat sejumlah anomali yang ditanam sengaja.
Pipeline memuatnya tanpa galat, dbt membangun tanpa galat, dan dashboard
menampilkan angka yang tampak wajar --- karena tidak satu pun anomali itu
melanggar constraint basis data.

\begin{center}
\small
\begin{tabular}{@{}p{0.22\textwidth}>{\raggedright\arraybackslash}p{0.34\textwidth}>{\raggedright\arraybackslash}p{0.34\textwidth}@{}}
\toprule
\textbf{Dimensi} & \textbf{Anomali yang ditanam} & \textbf{Mengapa lolos} \\
\midrule
Consistency & \perintah{subtotal} tidak sama dengan
  kuantitas $\times$ harga $-$ diskon &
  Ketiganya sendiri-sendiri masuk akal \\
Uniqueness & Satu struk terkirim dua kali dari sumber berbeda &
  Kunci gudang tetap unik \\
Validity & Harga satuan seribu kali lipat &
  Tetap bilangan positif \\
Timeliness & Sebagian transaksi sampai sepuluh hari terlambat &
  Tanggalnya sah \\
Completeness & Satu toko tidak mengirim data sama sekali &
  Ketiadaan tidak menghasilkan baris \\
\bottomrule
\end{tabular}
\end{center}

\begin{catatanasisten}
Daftar anomali di atas \emph{tidak} dibagikan ke mahasiswa. Menemukannya adalah
pekerjaan Bagian A, dan daftar ini dipakai asisten untuk memeriksa berapa yang
berhasil ditemukan.

Anomali completeness --- satu toko yang tidak mengirim apa pun --- adalah yang
paling jarang ditemukan mahasiswa, karena tidak ada baris yang bisa diperiksa.
Bila sampai menit ke-60 belum ada yang menemukannya, beri petunjuk berupa
pertanyaan: ``berapa toko yang mengirim data pada batch ini, dan berapa yang
mengirim pada batch sebelumnya?''
\end{catatanasisten}

\section{Alur Praktikum 120 Menit}

\begin{alurwaktu}
0--15 & Demonstrasi dashboard dengan angka yang tampak wajar tetapi salah. \\
15--95 & Kerja terbimbing: profiling, penyusunan uji, karantina, dan
  rekonsiliasi. \\
95--110 & Checkpoint otomatis dan pembacaan laporan kualitas. \\
110--120 & Penjelasan tugas scorecard kualitas data. \\
\end{alurwaktu}

\section{Prosedur Praktikum Terbimbing}

\subsection{Bagian A: Memprofilkan Batch Beranomali}

\langkah{A.1 Menjalankan profiling}{15 menit}

Jalankan \berkas{sql/modul-10/01-profiling.sql}, lalu bandingkan hasilnya
dengan profil batch sebelumnya. Yang dicari bukan hanya nilai yang mustahil,
melainkan juga nilai yang berubah tanpa sebab.

\begin{lstlisting}[style=sqlgaya]
-- Perbandingan antarbatch: yang berubah drastis adalah calon anomali.
SELECT batch_id,
       COUNT(*)                                  AS baris,
       COUNT(DISTINCT toko_id)                   AS toko_mengirim,
       ROUND(AVG(harga_satuan), 2)               AS rata_harga,
       ROUND(AVG(kuantitas), 3)                  AS rata_kuantitas,
       MAX(waktu_transaksi)::DATE                AS transaksi_terbaru
FROM   staging.penjualan_gabungan
GROUP  BY batch_id
ORDER  BY batch_id;
\end{lstlisting}

\langkah{A.2 Mencatat temuan}{10 menit}

Temukan sekurang-kurangnya \textbf{lima} anomali. Untuk masing-masing, catat
pada \berkas{temuan-anomali.md}: dimensi kualitas yang dilanggar, jumlah baris
yang terpengaruh, dan query yang membuktikannya.

\begin{catatan}
Satu di antara anomali yang ditanam tidak dapat ditemukan dengan memeriksa isi
tabel, karena barisnya memang tidak ada. Bila Anda hanya menemukan empat,
tanyakan pada diri sendiri: apa yang \emph{seharusnya} ada tetapi tidak ada?
\end{catatan}

\subsection{Bagian B: Menerjemahkan Aturan Bisnis Menjadi Uji}

\langkah{B.1 Menyusun daftar aturan}{10 menit}

Sebelum menulis kode, lengkapi tabel aturan pada
\berkas{aturan-kualitas.md}. Setiap baris memuat: aturan dalam bahasa bisnis,
dimensi kualitas, cara mengukurnya, ambang, dan tingkat keparahan beserta
alasannya.

\begin{center}
\small
\begin{tabular}{@{}p{0.28\textwidth}p{0.14\textwidth}p{0.12\textwidth}>{\raggedright\arraybackslash}p{0.34\textwidth}@{}}
\toprule
\textbf{Aturan} & \textbf{Dimensi} & \textbf{Tingkat} & \textbf{Alasan} \\
\midrule
Subtotal harus sama dengan hitungannya & Consistency & error &
  Angka laporan langsung salah \\
Satu produk hanya punya satu versi aktif & Uniqueness & error &
  Join dimensi menggandakan fakta \\
Baris unknown di bawah 5 persen & Completeness & warn &
  \ldots \\
\ldots & \ldots & \ldots & \ldots \\
\bottomrule
\end{tabular}
\end{center}

\langkah{B.2 Menulis empat generic test}{10 menit}

Keempatnya sudah Anda kenal dari Modul 8. Tambahkan sekarang tingkat
keparahannya.

\begin{lstlisting}[style=yamlgaya]
models:
  - name: fct_penjualan
    columns:
      - name: produk_key
        tests:
          - not_null:
              config:
                severity: error
          - relationships:
              to: ref('dim_produk')
              field: produk_key
      - name: status_batch
        tests:
          - accepted_values:
              values: ['lengkap', 'ada selisih', 'karantina']
              config:
                severity: warn
\end{lstlisting}

\langkah{B.3 Menulis empat custom test}{15 menit}

Uji singular adalah query yang mengembalikan baris \emph{yang melanggar}
aturan. Dua contoh disediakan pada \berkas{dbt/modul-10/tests/}; dua sisanya
Anda tulis sendiri.

\begin{lstlisting}[style=sqlgaya]
-- tests/satu_versi_scd_aktif.sql
-- Lulus bila query ini tidak mengembalikan baris apa pun.
SELECT produk_id, COUNT(*) AS jumlah_versi_aktif
FROM   {{ ref('dim_produk') }}
WHERE  baris_kini
GROUP  BY produk_id
HAVING COUNT(*) > 1
\end{lstlisting}

\begin{lstlisting}[style=sqlgaya]
-- tests/total_penjualan_konsisten.sql
-- Total mart harus sama dengan total fakta. Selisih berarti mart belum
-- disegarkan, atau agregasinya keliru.
WITH fakta AS (SELECT SUM(subtotal) AS n FROM {{ ref('fct_penjualan') }}),
     mart  AS (SELECT SUM(nilai_penjualan) AS n FROM {{ ref('agg_bulanan') }})
SELECT fakta.n AS total_fakta, mart.n AS total_mart
FROM   fakta, mart
WHERE  ROUND(fakta.n, 2) <> ROUND(mart.n, 2)
\end{lstlisting}

Dua uji yang Anda tulis sendiri harus mencakup dimensi yang belum tersentuh ---
misalnya \emph{consistency} untuk aritmetika subtotal, dan \emph{timeliness}
untuk selisih hari antara transaksi dan pemuatannya.

\begin{lstlisting}[style=shellgaya]
dbt test --store-failures
\end{lstlisting}

\begin{catatan}
Opsi \perintah{--store-failures} menyimpan baris yang gagal ke tabel
tersendiri, bukan hanya menghitungnya. Tanpa itu, Anda tahu ada 12 baris salah
tetapi tidak tahu baris yang mana --- dan tidak dapat memeriksanya.
\end{catatan}

\subsection{Bagian C: Mengarantina Data Gagal}

\langkah{C.1 Membuat tabel karantina}{10 menit}

\begin{lstlisting}[style=sqlgaya]
CREATE TABLE gudang.karantina_penjualan (
  karantina_id  BIGINT GENERATED ALWAYS AS IDENTITY PRIMARY KEY,
  batch_id      BIGINT      NOT NULL,
  aturan        VARCHAR(60) NOT NULL,   -- aturan yang dilanggar
  alasan        TEXT        NOT NULL,   -- keterangan bagi manusia
  ditemukan     TIMESTAMP   NOT NULL DEFAULT CURRENT_TIMESTAMP,
  ditinjau      BOOLEAN     NOT NULL DEFAULT FALSE,
  baris_asli    JSONB       NOT NULL    -- salinan utuh baris sumber
);
\end{lstlisting}

Kolom \perintah{baris\_asli} bertipe \perintah{JSONB} agar baris apa pun dapat
disimpan tanpa menuntut struktur yang sama. Yang disimpan adalah salinan ---
data sumber sendiri tidak disentuh.

\langkah{C.2 Mengarantina dan memuat ulang}{10 menit}

\begin{lstlisting}[style=sqlgaya]
-- Sisihkan baris yang melanggar aritmetika subtotal.
INSERT INTO gudang.karantina_penjualan (batch_id, aturan, alasan, baris_asli)
SELECT batch_id,
       'subtotal_konsisten',
       format('subtotal %s, seharusnya %s',
              subtotal, ROUND(kuantitas * harga_satuan - diskon, 2)),
       to_jsonb(s)
FROM   staging.penjualan_gabungan s
WHERE  ROUND(kuantitas * harga_satuan - diskon, 2) <> ROUND(subtotal, 2);

-- Muat ulang fakta, kecuali baris yang terkarantina.
DELETE FROM gudang.fakta_penjualan WHERE batch_id = :batch;
INSERT INTO gudang.fakta_penjualan (...)
SELECT ... FROM staging.penjualan_gabungan s
WHERE  s.batch_id = :batch
  AND  NOT EXISTS (SELECT 1 FROM gudang.karantina_penjualan k
                   WHERE k.batch_id = s.batch_id
                     AND k.baris_asli->>'transaksi_id'
                         = s.transaksi_id::TEXT);
\end{lstlisting}

Jumlah baris yang dikarantina menjadi salah satu penjelasan sah bagi selisih
rekonsiliasi pada Bagian D.

\subsection{Bagian D: Merekonsiliasi Sumber dan Gudang}

\langkah{D.1 Membandingkan empat besaran}{15 menit}

Rekonsiliasi lintas basis data kini mungkin berkat \perintah{postgres\_fdw}
yang Anda pasang pada Modul 7. Jalankan
\berkas{sql/modul-10/02-reconciliation.sql}.

\begin{lstlisting}[style=sqlgaya]
WITH sumber AS (
  SELECT COUNT(*)                       AS baris,
         COUNT(DISTINCT t.transaksi_id) AS transaksi_unik,
         SUM(ti.kuantitas)              AS kuantitas,
         SUM(ti.subtotal)               AS nilai
  FROM   src_oltp.transaksi_item ti
  JOIN   src_oltp.transaksi t ON t.transaksi_id = ti.transaksi_id
  WHERE  t.status <> 'BATAL'
    AND  t.waktu_transaksi::DATE = :tanggal_batch
),
gudang AS (
  SELECT COUNT(*)                     AS baris,
         COUNT(DISTINCT transaksi_id) AS transaksi_unik,
         SUM(kuantitas)               AS kuantitas,
         SUM(subtotal)                AS nilai
  FROM   gudang.fakta_penjualan
  WHERE  tanggal_key = :tanggal_key
)
SELECT 'baris'          AS besaran, s.baris,          g.baris,
       s.baris - g.baris                   AS selisih FROM sumber s, gudang g
UNION ALL
SELECT 'transaksi_unik', s.transaksi_unik, g.transaksi_unik,
       s.transaksi_unik - g.transaksi_unik FROM sumber s, gudang g
UNION ALL
SELECT 'kuantitas',      s.kuantitas,      g.kuantitas,
       s.kuantitas - g.kuantitas           FROM sumber s, gudang g
UNION ALL
SELECT 'nilai',          s.nilai,          g.nilai,
       s.nilai - g.nilai                   FROM sumber s, gudang g;
\end{lstlisting}

\langkah{D.2 Menjelaskan setiap selisih}{10 menit}

Untuk setiap selisih yang bukan nol, lengkapi tabel penjelasan pada
\berkas{laporan-kualitas.md}. Setiap baris harus menyebut penyebab dan
angkanya, dan jumlah seluruh penjelasan harus sama dengan selisihnya.

\begin{center}
\small
\begin{tabular}{@{}p{0.20\textwidth}r>{\raggedright\arraybackslash}p{0.48\textwidth}@{}}
\toprule
\textbf{Besaran} & \textbf{Selisih} & \textbf{Penjelasan} \\
\midrule
Baris & 47 & 39 baris terkarantina, 8 transaksi terlambat belum masuk \\
Nilai & \ldots & \ldots \\
\bottomrule
\end{tabular}
\end{center}

\begin{peringatan}
Selisih yang ``kira-kira sesuai'' tidak memenuhi syarat. Bila 47 baris tidak
cocok dan penjelasan Anda hanya mencakup 39, sisanya adalah kegagalan yang
belum ditemukan --- bukan pembulatan.
\end{peringatan}

\subsection{Bagian E: Menyusun Laporan Kualitas per Batch}

\langkah{E.1 Merangkum hasil}{10 menit}

Susun \berkas{laporan-kualitas.md} yang memuat, untuk batch ini: nilai keenam
dimensi kualitas beserta ambangnya, hasil seluruh uji, jumlah baris
terkarantina menurut aturan, dan tabel rekonsiliasi beserta penjelasannya.

\begin{lstlisting}[style=sqlgaya]
-- Scorecard satu batch, siap disalin ke laporan.
SELECT 'completeness' AS dimensi,
       ROUND(100.0 * COUNT(*) FILTER (WHERE produk_key = -1)
             / NULLIF(COUNT(*), 0), 3) AS nilai_persen,
       5.0 AS ambang_persen
FROM   gudang.fakta_penjualan
UNION ALL
SELECT 'consistency',
       ROUND(100.0 * COUNT(*) FILTER (
             WHERE ROUND(kuantitas * harga_satuan - diskon, 2)
                   <> ROUND(subtotal, 2)) / NULLIF(COUNT(*), 0), 3),
       0.0
FROM   gudang.fakta_penjualan;
\end{lstlisting}

\begin{luaran}
\berkas{modul-10-kualitas/temuan-anomali.md},
\berkas{aturan-kualitas.md}, proyek dbt berisi minimal delapan uji,
\berkas{ddl-karantina.sql}, dan \berkas{laporan-kualitas.md} beserta tabel
rekonsiliasi.
\end{luaran}

\begin{checkpoint}
Asisten menjalankan \berkas{sql/modul-10/check.sql} dan memeriksa butir berikut:
\begin{ceklis}
  \item seluruh anomali yang ditanam terdeteksi dan tercatat pada
        \berkas{temuan-anomali.md};
  \item sekurang-kurangnya delapan uji terpasang --- empat generic dan empat
        custom --- masing-masing dengan tingkat keparahan yang beralasan;
  \item seluruh uji bertingkat \perintah{error} lulus setelah remediasi;
  \item tabel karantina ada dan terisi, beserta aturan dan alasan setiap baris;
  \item keempat besaran rekonsiliasi diperiksa, dan setiap selisih dijelaskan
        sampai jumlahnya tepat;
  \item \textbf{data mentah tidak diubah} --- anomali bawaan masih utuh pada
        basis data sumber;
  \item laporan kualitas memuat keenam dimensi beserta ambangnya.
\end{ceklis}
\end{checkpoint}

\section{Latihan Mandiri di Kelas}

\latihan{Menakar ambang yang tepat}

Uji \emph{completeness} Anda melaporkan 3,2 persen baris fakta menunjuk baris
unknown.

\begin{enumerate}[leftmargin=1.4em]
  \item Pada ambang berapa Anda akan menghentikan pemuatan? Nyatakan alasannya
        dalam istilah dampak bisnis, bukan istilah teknis.
  \item Bila ambang ditetapkan 1 persen, berapa kali pipeline akan berhenti
        selama sebulan terakhir? Periksa dengan data batch yang ada.
  \item Apa yang terjadi pada disiplin tim bila pipeline berhenti terlalu
        sering? Kaitkan dengan pilihan antara \perintah{error} dan
        \perintah{warn}.
\end{enumerate}

\latihan{Membuat uji yang lolos padahal data salah}

Tulis satu uji yang \emph{lulus} untuk batch beranomali ini, padahal datanya
jelas bermasalah.

\begin{enumerate}[leftmargin=1.4em]
  \item Tunjukkan ujinya dan jelaskan mengapa ia lolos.
  \item Anomali mana yang luput darinya?
  \item Perbaiki ujinya sehingga anomali itu tertangkap.
  \item Pelajaran apa yang Anda tarik tentang uji yang selalu lulus?
\end{enumerate}

\section{Tugas Individu}

\subsection{Pekerjaan Wajib}

Mahasiswa menyusun scorecard kualitas untuk proses bisnis kedua, menetapkan
threshold beserta alasannya, dan membandingkan hasil pengukuran pada dua batch.

Kumpulkan \berkas{scorecard-kualitas.md} yang memuat:

\begin{enumerate}[leftmargin=1.4em]
  \item proses bisnis kedua yang dipilih --- persediaan, promosi, atau retur;
  \item keenam dimensi kualitas beserta cara mengukurnya pada proses tersebut,
        lengkap dengan query masing-masing;
  \item ambang setiap dimensi beserta alasannya dalam istilah dampak bisnis,
        dan tingkat keparahan yang dipilih;
  \item hasil pengukuran pada \textbf{dua} batch berbeda, disandingkan dalam
        satu tabel;
  \item penilaian apakah keadaan membaik atau memburuk, beserta satu tindakan
        yang Anda usulkan.
\end{enumerate}

\begin{catatan}
Butir keempat adalah alasan scorecard dibuat. Satu batch hanya memberi potret;
dua batch memberi arah. Sebagian besar keputusan kualitas data diambil
berdasarkan arah, bukan potret --- ``tiga persen'' tidak berarti apa-apa
sampai diketahui bulan lalu berapa.
\end{catatan}

\subsection{Pertanyaan Analisis}

\begin{enumerate}[leftmargin=1.4em]
  \item Anomali yang paling sulit Anda temukan adalah toko yang tidak mengirim
        data sama sekali. Jelaskan mengapa jenis kesalahan ini luput dari
        seluruh uji atas isi tabel, dan rancang satu pemeriksaan yang akan
        menangkapnya secara otomatis pada batch berikutnya.
  \item Seorang rekan mengusulkan memperbaiki 39 baris bermasalah langsung di
        basis data sumber agar rekonsiliasi cocok sempurna. Uraikan apa yang
        hilang bila usulan itu dijalankan, dan sebutkan pihak mana yang
        seharusnya berwenang memperbaiki data sumber.
  \item Seluruh uji Anda lulus pada batch berikutnya. Sebutkan dua sebab yang
        mungkin selain ``datanya memang bersih'', dan pemeriksaan apa yang akan
        Anda lakukan untuk memastikannya.
\end{enumerate}

\section{Format Pengumpulan dan Checklist}

Kumpulkan skrip profiling dan rekonsiliasi, proyek dbt berisi seluruh uji, DDL
karantina, dan laporan kualitas per batch.

Seluruh berkas diletakkan pada \berkas{modul-10-kualitas/}.

\begin{ceklis}
  \item \berkas{pre-lab.md}
  \item \berkas{temuan-anomali.md} --- lima anomali beserta bukti angka
  \item \berkas{aturan-kualitas.md} --- aturan, dimensi, ambang, keparahan,
        alasan
  \item Proyek dbt berisi minimal delapan uji
  \item \berkas{ddl-karantina.sql} dan isi tabel karantina
  \item \berkas{laporan-kualitas.md} --- scorecard dan tabel rekonsiliasi
  \item \berkas{scorecard-kualitas.md} --- proses bisnis kedua, dua batch
  \item \berkas{jawaban-analisis.md}
  \item Bukti bahwa data sumber tidak berubah
\end{ceklis}

\section{Troubleshooting}

\begin{table}[htbp]
\centering
\small
\caption{Masalah yang lazim muncul pada Modul 10 dan penanganannya}
\label{tab:m10-troubleshooting}
\begin{tabular}{@{}p{0.30\textwidth}>{\raggedright\arraybackslash}p{0.62\textwidth}@{}}
\toprule
\textbf{Gejala} & \textbf{Penanganan} \\
\midrule
Seluruh uji lulus pada batch beranomali &
  Uji terlalu longgar, atau menguji hal yang memang tidak dilanggar. Bandingkan
  daftar uji Anda dengan keenam dimensi --- biasanya ada dimensi yang belum
  tersentuh sama sekali. \\
Uji singular selalu gagal padahal data benar &
  Logikanya terbalik. Uji singular mengembalikan baris yang \emph{melanggar};
  bila query Anda mengembalikan baris yang benar, ia akan selalu gagal. \\
\perintah{dbt test} melaporkan kegagalan tanpa menyebut barisnya &
  Jalankan dengan \perintah{--store-failures}, lalu periksa tabel hasilnya. \\
Rekonsiliasi selisihnya nol padahal ada anomali &
  Anomali yang ditanam tidak mengubah jumlah maupun total --- misalnya baris
  yang tertukar. Karena itu keempat besaran diperiksa, bukan satu. \\
Selisih rekonsiliasi tidak dapat dijelaskan seluruhnya &
  Periksa berurutan: baris terkarantina, transaksi batal, transaksi terlambat,
  lalu baris yang menunjuk unknown. \\
Rekonsiliasi lambat sekali &
  Query menempuh \perintah{postgres\_fdw} dan tidak seluruh agregasi dapat
  didorong ke sisi seberang. Batasi pada satu tanggal batch, jangan seluruh
  riwayat. \\
Jumlah baris karantina bertambah setiap pemuatan diulang &
  Penyisihan tidak idempoten. Hapus baris karantina batch itu lebih dahulu,
  sebagaimana pola hapus-lalu-muat Modul 7. \\
Anomali bawaan hilang dari sumber &
  Seseorang memperbaikinya langsung di sumber. Ini melanggar aturan modul;
  pulihkan sumber dari seed, lalu ulangi dari Bagian A. \\
\bottomrule
\end{tabular}
\end{table}

\section{Ringkasan}

Modul ini menutup rangkaian sepuluh modul dengan membalik andaian yang berlaku
sejak awal: data tidak dianggap benar sampai dibuktikan. Pembuktian itu berupa
angka, dijalankan otomatis, dan diulang setiap kali data baru masuk.

Tiga gagasan menutup praktikum ini. \emph{Pertama}, ``berkualitas'' harus
dipecah menjadi dimensi yang masing-masing terukur --- karena hanya angka yang
memungkinkan penetapan ambang, perbandingan antarbatch, dan penilaian arah.
\emph{Kedua}, kesalahan yang paling berbahaya adalah yang tidak melanggar satu
constraint pun: subtotal yang tidak konsisten, harga yang seribu kali lipat,
dan toko yang diam. Ketiganya lolos sampai ke dashboard tanpa satu pun galat.
\emph{Ketiga}, data mentah tidak diubah --- baris gagal dikarantina, bukan
dibuang atau diperbaiki di sumber, karena rekonsiliasi yang cocok akibat sumber
yang disunting tidak membuktikan apa pun.

Sepuluh modul ini membangun satu gudang data secara berlapis: dari membaca
sumber pada Modul 1, sampai membuktikan bahwa angkanya layak dipercaya hari
ini. Proyek akhir pada minggu 14 dan 15 mematangkan apa yang sudah ada ---
menerapkannya pada proses bisnis kedua, melengkapinya dengan runbook
operasional, dan membuktikan satu siklus muat berulang yang idempoten.
